\documentclass[11pt]{article}
\usepackage[utf8]{inputenc}  
\usepackage[dutch]{babel}    
\usepackage{hyperref}        
\usepackage{booktabs}		 
\usepackage{a4wide}          
\pagestyle{empty}            
\usepackage{graphicx}        
\usepackage{grffile}         
\usepackage{caption}         
\usepackage{float}
\usepackage{fancyhdr}
\usepackage{titling}
\setlength{\droptitle}{-2cm} 
\posttitle{\end{center}\vspace{-0.5cm}} 
\postauthor{\end{center}\vspace{-0.5cm}} 
\predate{\begin{center}}
\postdate{\end{center}\vspace{-0.8cm}} 

\title{Digital Electronic Systems Design Project Report}
\author{Federico Ferro, Simone Guan, Xiao Chen Wang, Cheng Marco Zhang}
\date{Academic Year 2025/2026} 

\begin{document}

\fancypagestyle{plain}{
    \fancyhf{}
    \lfoot{\footnotesize Politecnico di Milano}
    \rfoot{\footnotesize Spring 2026: Digital Electronic Systems Design}
    \renewcommand{\headrulewidth}{0pt}
    \renewcommand{\footrulewidth}{0pt}
}


\maketitle               
\section*{Timing Analysis}

% test
From our implementation, the timing analysis yields compliant results for both setup and hold constraints at the target operating frequency of 150~MHz.

Regarding the setup time, the most critical path is located within the moving average block (Path 21). While the design requirement demands a 150~MHz clock, this implementation achieves a positive slack of 0.190~ns. This safety margin allows the circuit to potentially operate at a maximum frequency of approximately 154~MHz.

Concerning the hold time, the critical path is identified along the data route from the \texttt{diligent\_jstk2} module to the \texttt{effect\_selector\_0} unit (Path 31). The required hold time constraint is 0~ns, and our implementation guarantees a safe positive hold slack of 0.026~ns, ensuring that no race conditions occur.

The detailed parameters resulting from the implementation are summarized in Table 1 for setup time and Table 2 for hold time below.
\begin{table}[H]
    \centering 
    \begin{minipage}[b]{0.48\textwidth}
        \flushleft
        \captionsetup{type=table, singlelinecheck=false, justification=raggedright, skip=2pt}
        \caption{Timing Summary: Setup Time (Path 21)}
        \label{tab:setup_report}
        \small
        \begin{tabular}{lrc}
            \toprule
            \textbf{Parameter} & \textbf{Value} & \textbf{Unit} \\
            \midrule
            Slack             & 0.190 & ns \\
            Requirement       & 6.731 & ns \\
            Data Path Delay   & 6.519 & ns \\
            \quad \textit{Logic Delay} & 2.916 & ns \\
            \quad \textit{Route Delay} & 3.603 & ns \\
            Clock Path Skew   & 0.063 & ns \\
            Clock Uncertainty & 0.116 & ns \\
            \bottomrule
        \end{tabular}
    \end{minipage}%
    \hfill 
    \begin{minipage}[b]{0.48\textwidth}
        \flushleft
        \captionsetup{type=table, singlelinecheck=false, justification=raggedright, skip=2pt}
        \caption{Timing Summary: Hold Time (Path 31)}
        \label{tab:hold_report}
        \small
        \begin{tabular}{lrc}
            \toprule
            \textbf{Parameter} & \textbf{Value} & \textbf{Unit} \\
            \midrule
            Slack             & 0.026 & ns \\
            Requirement       & 0.000 & ns \\
            Data Path Delay   & 0.360 & ns \\
            \quad \textit{Logic Delay} & 0.141 & ns \\
            \quad \textit{Route Delay} & 0.219 & ns \\
            Clock Path Skew   & 0.263 & ns \\
            Clock Uncertainty & 0.116 & ns \\ 
            \bottomrule
        \end{tabular}
    \end{minipage}
    
\end{table}


\section*{Power Analysis}

At the operating frequency of 150 MHz, the following implementation's power consumption is as showed in the following tabels.

\begin{table}[H]
    \centering 
    \begin{minipage}[b]{0.46\textwidth}
        \flushleft
        \captionsetup{type=table, singlelinecheck=false, justification=raggedright, skip=2pt}
        \caption{Power Analysis Summary.}
        \label{tab:main_power}
        \small
        \begin{tabular}{lrc}
            \toprule
            \textbf{Parameter} & \textbf{Value} & \textbf{Unit} \\
            \midrule
            Total On-Chip Power            & 0.325 & W   \\
            Design Power Budget            & N/A   & --  \\
            Power Budget Margin            & N/A   & --  \\
            Junction Temperature           & 26.6  & °C  \\
            Thermal Margin                 & 58.4  & °C  \\
            Effective $\theta_{JA}$        & 5.0   & °C/W\\
            Power to Off-Chip Devices      & 0     & W   \\
            Confidence Level               & Low   & --  \\
            \bottomrule
        \end{tabular}
    \end{minipage}%
    \hfill 
    \begin{minipage}[b]{0.50\textwidth}
        \flushleft
        \captionsetup{type=table, singlelinecheck=false, justification=raggedright, skip=2pt}
        \caption{On-Chip Power Breakdown.}
        \label{tab:power_dynamic_breakdown}
        \small
        \begin{tabular}{lrcc}
            \toprule
            \textbf{Component} & \textbf{Power} & \textbf{Unit} & \textbf{Ratio (\%)} \\
            \midrule
            \textbf{Dynamic}   & \textbf{0.252} & \textbf{W} & \textbf{77\%} \\
            \quad \textit{MMCM}               & 0.125 & W & 50\% \\
            \quad \textit{BRAM}               & 0.077 & W & 30\% \\
            \quad \textit{Clocks}             & 0.017 & W &  7\% \\
            \quad \textit{I/O}                & 0.016 & W &  6\% \\
            \quad \textit{Signals}            & 0.010 & W &  4\% \\
            \quad \textit{Logic}              & 0.008 & W &  3\% \\
            \quad \textit{DSP}                & 0.000 & W &  0\% \\
            \midrule
            \textbf{Device Static} & \textbf{0.074} & \textbf{W} & \textbf{23\%} \\
            \bottomrule
        \end{tabular}
    \end{minipage}
    
\end{table}

\end{document}
